\Needspace{5\baselineskip}\textbf{Problems}

\begin{enumerate}[label=\thesection.\arabic*.,ref=\thesection.\arabic*,start=5,series=book-problems]

\item\label{ex:mean-zero-neumann-variational-compatibility} 假设 $\Omega$ 有界, 连通且具有 $1$ 阶延拓性质. 令\[V=\left\{u\in H^1(\Omega)\ \middle|\ \int_\Omega u(\bm{x})\,\mathrm{d}\bm{x}=0\right\}.\](a) 证明, 对每个 $f\in L^2(\Omega)$, 都存在唯一的 $u\in V$, 使得
\begin{equation}\label{eq:mean-zero-neumann-variational-problem}
\int_\Omega\nabla u\cdot\nabla v=\int_\Omega fv\quad\text{for all }v\in V.
\end{equation}
(见 Problem~\ref{ex:mean-zero-sobolev-poincare-inequality}.)

(b) 说明当 $\int_\Omega f=0$ 时, 为什么适合将 \eqref{eq:mean-zero-neumann-variational-problem} 看作 Neumann 问题
\begin{equation}\label{eq:mean-zero-neumann-strong-problem}
\begin{aligned}
-\Delta u&=f&&\text{in }\Omega\\
\frac{\partial u}{\partial n}&=0&&\text{on }\partial\Omega
\end{aligned}
\end{equation}
的弱形式.

(c) 若 $\int_\Omega f\ne0$, 仍将 \eqref{eq:mean-zero-neumann-variational-problem} 的解称为 \eqref{eq:mean-zero-neumann-strong-problem} 的解是否合理? 请说明.

\item\label{ex:scalar-dirichlet-complementing-condition} 证明, 标量椭圆 PDE 的 Dirichlet 边界条件总满足补条件.

\item\label{ex:dirichlet-system-triangular-reduction-details} 补全 Lemma~\ref{lem:dirichlet-system-tangential-triangular-reduction} 证明的细节.

\item\label{ex:agmon-condition-implies-complementing-natural-boundaries} 假设 Agmon 条件成立. 证明 \eqref{eq:agmon-essential-natural-boundary-value-problem} 满足补条件. Hint: 在半空间上应用 \eqref{eq:agmon-constrained-form-green-boundary-identity}.

\item\label{ex:inhomogeneous-essential-natural-boundaries-weak-form} 允许边界条件非齐次时, 给出 \eqref{eq:agmon-essential-natural-boundary-value-problem} 的弱形式.

\item\label{ex:stokes-traction-complementing-condition} 证明, ``牵引边界条件\textquotesingle\textquotesingle\ $(\nabla\bm{u}+(\nabla\bm{u})^{\mathrm{T}})\cdot\bm{n}-p\bm{n}=0$ 对 Stokes 方程组满足补条件.

\item\label{ex:scalar-dirichlet-zero-fredholm-index} 证明, 具有 Dirichlet 条件的标量椭圆算子的 Fredholm 指标为 $0$. Hint: 证明伴随问题也具有 Dirichlet 条件.

\end{enumerate}
